\begin{summarybox}
	\textbf{\textcolor{CSummary}{一句话核心}}
	通过拆项、添项或配系数，使变形后的正项乘积为常数，再利用均值不等式即可得到和的最小值；等号条件为各项相等。

	\medskip
	\textbf{\textcolor{CSummary}{使用条件}}
	\begin{itemize}[leftmargin=*, itemsep=0.5em, topsep=0.4em]
		\item \textbf{条件1（各项正数）：} 变形后的每一项都必须大于零。
		\item \textbf{条件2（乘积常数）：} 所有正项的乘积必须是常数（与变量无关）。
		\item \textbf{条件3（等号可达）：} 在变量允许取值范围内，能够使各项相等（否则需谨慎使用）。
	\end{itemize}

	\medskip
	\textbf{\textcolor{CSummary}{关键提醒}}
	\begin{itemize}[leftmargin=*, itemsep=0.5em, topsep=0.4em]
		\item \textbf{易错点（正负与乘积）：} 注意检查各项正负和乘积是否为定值，二者缺一不可。
		\item \textbf{检查项（取等验证）：} 使用后必须验证等号能否在定义域内取到，避免最小值虚设。
	\end{itemize}

\end{summarybox}
